\section{EDT}
La EDT organiza el alcance completo en entregables con una sola imputación de esfuerzo. Esta sección presenta las ramas principales y un ejemplo de descomposición; el Formulario T-14 desarrolla resultados, aceptación y responsables por paquete. La trazabilidad enlaza requisitos, construcción, pruebas e incorporación al servicio.

\subsection{Criterio de descomposición y código de cuentas}
Synaptix adopta una descomposición por entregables. El código raíz \texttt{1} identifica el contrato; los códigos \texttt{1.n} agrupan entregables principales; y los niveles inferiores identifican componentes y paquetes. La etapa contractual será un atributo del trabajo, de modo que una capacidad desarrollada progresivamente conserve su relación con el mismo entregable. La EDT expresa pertenencia y cobertura; las fechas y dependencias se registran en la red de actividades.

La descomposición se detendrá cuando el paquete tenga un resultado medible, una condición de aceptación, un responsable único y una estimación defendible. Las instancias actualmente estimadas tienen entre 8 y 80 horas-persona; este rango describe la estimación y no constituye una exigencia de las Bases ni un límite para definir nuevos paquetes. Se mantiene un período quincenal de control. La dedicación parcial y las esperas externas se declararán por separado; el esfuerzo del paquete no equivale automáticamente a su duración. Los servicios recurrentes se presentan mediante familias con resultado, responsable, frecuencia y aceptación comunes; cada período conserva su instancia estimada y asignada.

Cada rama incluirá todo el trabajo de su componente superior y excluirá el trabajo imputado a otra rama. La verificación de cobertura cruzará requisitos, entregables y paquetes; la suma de porcentajes de un gráfico no sustituirá esa comprobación. Las pruebas integradas, los servicios compartidos y las innovaciones tendrán una imputación única; el uso de un componente por varios dominios se registrará como dependencia.

\subsection{Inventario de entregables principales}
El inventario organiza los productos y servicios de SD3, SD4, SD5 y SD13. La Figura~\ref{fig:sd7-edt-general} presenta las doce ramas y la Tabla~\ref{tab:sd7-inventario} relaciona cada cuenta con su fuente de alcance. Estos códigos agrupan entregables; los paquetes terminales se identifican en T-14.

\begin{figure}[htbp]
\centering
\begin{tikzpicture}[font=\fontsize{9}{11}\selectfont,
 caja/.style={draw=SynLine,fill=SynSurface,text width=42mm,minimum height=12mm,align=center,inner sep=2mm}]
\node[caja,fill=SynNavy,text=SynWhite,text width=132mm] (raiz) at (0,1.5) {\textbf{1 · Contrato de solución y servicio de la marina}};
\node[caja] (c0) at (-5.0,0.0) {\textbf{1.1}\\Dirección y control};
\node[caja] (c1) at (0,0.0) {\textbf{1.2}\\Plataforma híbrida};
\node[caja] (c2) at (5.0,0.0) {\textbf{1.3}\\Auditoría y seguridad};
\node[caja] (c3) at (-5.0,-1.6) {\textbf{1.4}\\Núcleo Operacional};
\node[caja] (c4) at (0,-1.6) {\textbf{1.5}\\Núcleo de Servicios};
\node[caja] (c5) at (5.0,-1.6) {\textbf{1.6}\\Núcleo Administrativo};
\node[caja] (c6) at (-5.0,-3.2) {\textbf{1.7}\\Canales e integración};
\node[caja] (c7) at (0,-3.2) {\textbf{1.8}\\Migración e historia};
\node[caja] (c8) at (5.0,-3.2) {\textbf{1.9}\\Innovaciones};
\node[caja] (c9) at (-5.0,-4.800000000000001) {\textbf{1.10}\\Verificación y aceptación};
\node[caja] (c10) at (0,-4.800000000000001) {\textbf{1.11}\\Formación e implantación};
\node[caja] (c11) at (5.0,-4.800000000000001) {\textbf{1.12}\\Servicio y salida};
\draw[SynBlue] (raiz.south) -- (0,.8);
\draw[SynBlue] (-7.5,.8) -- (2.5,.8);
\draw[SynBlue] (-7.5,.8) -- (-7.5,-4.8);
\draw[SynBlue] (-7.5,0.0) -- (c0.west);
\draw[SynBlue] (-7.5,-1.6) -- (c3.west);
\draw[SynBlue] (-7.5,-3.2) -- (c6.west);
\draw[SynBlue] (-7.5,-4.8) -- (c9.west);
\draw[SynBlue] (-2.5,.8) -- (-2.5,-4.8);
\draw[SynBlue] (-2.5,0.0) -- (c1.west);
\draw[SynBlue] (-2.5,-1.6) -- (c4.west);
\draw[SynBlue] (-2.5,-3.2) -- (c7.west);
\draw[SynBlue] (-2.5,-4.8) -- (c10.west);
\draw[SynBlue] (2.5,.8) -- (2.5,-4.8);
\draw[SynBlue] (2.5,0.0) -- (c2.west);
\draw[SynBlue] (2.5,-1.6) -- (c5.west);
\draw[SynBlue] (2.5,-3.2) -- (c8.west);
\draw[SynBlue] (2.5,-4.8) -- (c11.west);
\end{tikzpicture}
\caption{Vista general de las doce cuentas de entregables de la EDT}
\label{fig:sd7-edt-general}
\FuenteFigura{Elaboración propia de Synaptix a partir del catálogo integrado de T-14.}
\end{figure}

Las doce cajas pertenecen directamente a la raíz 1; las líneas enlazan cada cuenta a esa raíz y su posición en columnas no establece precedencias entre cuentas. Cada rama se descompone hasta sus paquetes en T-14; 7.3 desarrolla la red temporal.

\begin{table}[H]
\centering
\begin{tabularx}{\linewidth}{L{12mm}Y L{31mm}}
\hline
 \EncabezadoTabla{Cuenta} & \EncabezadoTabla{Entregable} & \EncabezadoTabla{Origen del alcance}\\\hline
1.1 & Expediente de dirección y control & SD6; T-9\\\hline
1.2 & Plataforma híbrida habilitada & SD4; T-11\\\hline
1.3 & Controles de auditoría y seguridad habilitados & SD4; SD5\\\hline
1.4 & Capacidades del Núcleo Operacional & SD3; SD5\\\hline
1.5 & Capacidades del Núcleo de Servicios & SD3; SD5\\\hline
1.6 & Capacidades del Núcleo Administrativo & SD3; SD5\\\hline
1.7 & Canales y servicios comunes integrados & SD3; SD4\\\hline
1.8 & Datos migrados y archivo histórico consultable & SD5\\\hline
1.9 & Cinco innovaciones incorporadas y evaluadas & SD13; T-19\\\hline
1.10 & Expediente de verificación y aceptación & SD6; SD9; T-17\\\hline
1.11 & Usuarios preparados y etapas implantadas & SD3; T-18\\\hline
1.12 & Servicio gestionado y transferencia de salida & SD4; SD5; T-15\\\hline
\end{tabularx}
\caption{Correspondencia de cuentas y entregables con su alcance de origen}
\label{tab:sd7-inventario}
\FuenteTabla{Elaboración propia de Synaptix a partir de SD3--SD6, SD9, SD13 y los formularios indicados.}
\end{table}

La correspondencia permite localizar la autoridad de cada entregable y evita estimar dos veces una capacidad compartida. La cuenta 1.1 reúne línea base, trazabilidad, actas, cambios, riesgos, informes y control físico y financiero, incluido el plan de reversibilidad de los primeros 90 días. La plataforma 1.2 comprende DEV, QA, PREPROD, PROD y DR, nube, recinto, terminales, conectividad, continuidad y observabilidad; SD4/T-11 delimitan los activos. La cuenta 1.3 desarrolla identidades, permisos, consentimiento, cifrado, claves y auditoría; sus ensayos integrales se imputan a 1.10.

El Núcleo Operacional integra Operación marítima, Escuela náutica e Infraestructura: salidas y regresos, amarras, participación y retiro, inventario físico e inspección de boyas sin conexión. El Núcleo de Servicios incorpora Varadero y contratistas, Combustible y Ambiente y consumos; la medición inicial corresponde a E1 y las capacidades completas a E2. El Núcleo Administrativo comprende Clientes y contratos y Finanzas y cobranza, preservando la autoridad tributaria del sistema contable externo.

Los canales Android, portal y consola, interfaces externas, notificaciones, telemetría y sincronización se construyen una vez en 1.7; sus reglas por dominio permanecen en 1.4--1.6. La cuenta 1.8 comprende diagnóstico, herramientas, saneamiento, cargas, conciliación y consulta independiente del sistema de 2014, con carga inicial E1 y cobertura histórica E2 según SD5.

Las cinco innovaciones se incorporan en 1.9 con su esfuerzo incremental: solicitudes de mantenimiento, mantenimiento preventivo en invernada, lectura asistida, servicio gestionado por resultados y Pasaporte Ambiental Náutico. La cuenta 1.10 reúne protocolos, ensayos entre componentes y evidencia para T-17, incluidos los dos ensayos completos de migración de SD5. Formación, olas, convivencia, conciliación diaria, marcha blanca, reversión y estabilización se imputan a 1.11. El servicio 1.12 conserva soporte, NOC/SOC, mantenimiento, respaldos, SLA y transferencia final durante 21--56.

El inventario separa la construcción de capacidades de su comprobación integral. Cada paquete de construcción incluirá sus pruebas unitarias y documentación; 1.10 reunirá los ensayos entre componentes y la evidencia de recepción. Las pruebas internas específicas de una innovación permanecerán en 1.9 y sus resultados alimentarán el expediente general, sin volver a estimar la misma ejecución.

Las compras y obras físicas atribuidas al cliente por el Subdocumento 2 se registrarán como habilitaciones externas con responsable, fecha requerida y condiciones de recepción. Synaptix mantendrá en sus paquetes las especificaciones, la coordinación y la verificación de integración que le corresponden. Estas dependencias no se tratarán como instalaciones ya ejecutadas ni como ampliaciones automáticas del plazo \parencite[cap.~11, p.~27; cap.~17, secc.~17.5, p.~42]{caso07}.

Los códigos \texttt{IN1} a \texttt{IN5} del Subdocumento 13 conservarán su identidad como referencias de origen. El trabajo de diseño, construcción, pruebas y pilotos de las innovaciones se imputará a 1.9; el servicio recurrente, a 1.12. En particular, los informes \texttt{IN1.7.m.q} y la operación \texttt{IN2.7} y \texttt{IN3.7} conservarán su referencia a SD13 dentro de la rama de servicio. Cada paquete de ejecución tendrá una sola imputación de esfuerzo; si un componente de SD13 requiere subdivisión, la matriz de correspondencia identificará todos sus paquetes hijos sin contar también el esfuerzo del padre.

El calendario de la innovación 2 conservará $M_{\mathrm{inv}}$ como el primer junio disponible con IN1 y sus interfaces aceptadas, no anterior al mes 21, y mantendrá la secuencia de invernada y botaduras definida en SD13. Los componentes \texttt{IN4.4} e \texttt{IN5.4} distinguirán sus resultados de puesta en servicio de sus trabajos recurrentes. T-14 conserva la correspondencia detallada de los códigos y T-15 sus cargas y asignaciones; los padres no reciben una segunda imputación.

\subsection{Expediente de diagnóstico de las fuentes de migración}
El componente \texttt{1.8.1} entregará un expediente que permita decidir cómo extraer, interpretar y conciliar los antecedentes antes de construir las cargas definitivas. Su cobertura comprende los seis conjuntos obligatorios de SD5, sección 5.3.1: socios y armadores; contratos vigentes y documentos; seis años de pagos y saldos; embarcaciones e invernadas; los 14 expedientes de abandono; y tres años de matrícula escolar. El inventario abarcará las fuentes de los doce años de historia para delimitar su tratamiento, sin extender automáticamente a todo ese período la carga operacional \parencite[cap.~15, RT-05.15, p.~34]{caso07}.

La Figura~\ref{fig:sd7-rama-diagnostico} muestra la pertenencia de los cuatro paquetes al componente 1.8.1. La Tabla~\ref{tab:sd7-paquetes-diagnostico} resume sus resultados y esfuerzos. Las horas son una estimación ascendente inicial de la oferta, basada en las tareas enumeradas en el diccionario; no son mediciones de productividad ni duración comprometida.

\begin{figure}[htbp]
\centering
\begin{tikzpicture}[
 paquete/.style={draw=SynLine,fill=SynSurface,rounded corners=1mm,
 text width=82mm,minimum height=13mm,inner sep=3mm,align=left,
 font=\fontsize{9}{12}\selectfont},
 rama/.style={draw=SynNavy,fill=SynNavy,text=SynWhite,
 text width=32mm,inner sep=3mm,align=left,font=\fontsize{9}{12}\selectfont}]
\node[rama] (raiz) at (-5,0) {\textbf{1.8.1}\\Expediente de diagnóstico de las fuentes de migración};
\node[paquete] (p1) at (2,2.7) {\textbf{1.8.1.1}\\Catálogo de fuentes y accesos};
\node[paquete] (p2) at (2,.9) {\textbf{1.8.1.2}\\Procedimiento de extracción diagnóstica};
\node[paquete] (p3) at (2,-.9) {\textbf{1.8.1.3}\\Informe reproducible de calidad del origen};
\node[paquete] (p4) at (2,-2.7) {\textbf{1.8.1.4}\\Matriz de cobertura y decisiones de migración};
\draw[SynBlue,line width=.8pt] (raiz.east) -- (-2.6,0);
\draw[SynBlue,line width=.8pt] (-2.6,2.7) -- (-2.6,-2.7);
\foreach \paquete/\altura in {p1/2.7,p2/.9,p3/-.9,p4/-2.7}{
 \draw[SynBlue,line width=.8pt] (-2.6,\altura) -- (\paquete.west);
}
\end{tikzpicture}
\caption{Descomposición del expediente de diagnóstico en paquetes}
\label{fig:sd7-rama-diagnostico}
\FuenteFigura{EDT propuesta por Synaptix para el diagnóstico previo a la migración.}
\end{figure}


La jerarquía asigna cada paquete a un único componente superior. Las líneas expresan esa pertenencia; las precedencias de ejecución se conservan en la red de actividades de T-14/T-15.

\begin{table}[H]
\centering
\begin{tabularx}{\linewidth}{L{25mm}YR{13mm}L{31mm}}
\hline
 \EncabezadoTabla{Código} & \EncabezadoTabla{Resultado} & \EncabezadoTabla{Horas} & \EncabezadoTabla{Responsable}\\\hline
1.8.1.1 & Catálogo de fuentes y accesos & 24 & Líder de Datos\\\hline
1.8.1.2 & Procedimiento de extracción diagnóstica & 40 & Líder de Datos\\\hline
1.8.1.3 & Informe reproducible de calidad del origen & 32 & Líder de Datos\\\hline
1.8.1.4 & Matriz de cobertura y decisiones de migración & 24 & Líder de Datos\\\hline
\textbf{Total} & \textbf{Expediente de diagnóstico} & \textbf{120} & \\\hline
\end{tabularx}
\caption{Paquetes del expediente de diagnóstico de fuentes}
\label{tab:sd7-paquetes-diagnostico}
\FuenteTabla{Estimación ascendente propuesta por Synaptix; alcance definido en SD5, sección 5.3.}
\end{table}

Las 120 horas comprenden 24 de Análisis Funcional, 72 de Datos, 8 de Seguridad y 16 de Calidad. Incluyen documentación y comprobación de los productos de diagnóstico. No incluyen saneamiento masivo, carga definitiva, consulta histórica implementada, los dos ensayos completos de migración ni la operación del servicio; esos trabajos requieren otros componentes de 1.8, 1.10 y 1.12. La aceptación de este expediente no equivale a la aceptación de la migración.

\ifdefined\EnFormularioTCatorce
El diccionario detallado de los cuatro paquetes se presenta en el apartado siguiente de este formulario.
\else
El diccionario detallado de los cuatro paquetes se presenta en el archivo independiente \path{SYNAPTIX-Formulario-T-14.pdf}.
\fi
Sus fichas conservan estos códigos, esfuerzos, responsables y fronteras; el resumen no sustituye sus condiciones de aceptación.

El catálogo consolidado contiene 2.181 paquetes terminales, cada uno con resultado, aceptación, responsable único y 8--80 HH. El diccionario integrado enlaza las fichas específicas y las instanciaciones periódicas. La trazabilidad incluye los 576 IDs de T-12, los 27 complementos C07 y las 29 exclusiones Deseables; ninguna obligación de base se reduce por prioridad funcional. Cada prueba recibe el enunciado completo, sus cláusulas y el componente ofrecido. La existencia del cruce no se presenta como prueba ya ejecutada ni modifica las declaraciones de T-12.

Los padres no reciben horas, y los paquetes compartidos se imputan una vez. Los nueve dominios mantienen sus autoridades de SD3/SD4/SD5; las innovaciones aportan capacidades incrementales. Plataforma incluye DEV, QA, PREPROD, PROD y DR; la medición individual y su recepción se anticipan para disponer de serie previa a 2029. Las obligaciones corporativas de oferta se distinguen del alcance del contrato.
